\SourcePage{340}{343}
\section{X-ray terms}

An atom's inner electrons have such large binding energies that, when one of them passes into an unfilled outer shell (or is removed from the atom altogether), the resulting excited atom (or ion) is mechanically unstable against ionization accompanied by rearrangement of the electron shell and formation of a stable ion. Electron interactions within an atom are comparatively weak, however, so the probability of such a transition is still rather small and the excited state has a long lifetime $\tau$. Its level ``width'' $\hbar/\tau$ (see \S~44) is therefore small enough that the energies of an atom with an excited inner electron can meaningfully be treated as discrete energy levels of ``quasistationary'' atomic states. These levels are called \emph{X-ray terms}\footnote{The name reflects the fact that transitions between these levels cause the atom to emit X-rays.}.

X-ray terms are classified first by specifying the shell from which the electron has been removed, or, as one says, the shell in which a hole has formed. The particular destination of the electron has almost no effect on the atomic energy and is consequently immaterial.

The total angular momentum of all the electrons filling a shell is zero. Removing one electron gives the shell a certain angular momentum $j$. For the shell $(n,l)$, $j$ may take the values $l\pm1/2$. The resulting levels could thus be denoted by $1s_{1/2}$, $2s_{1/2}$, $2p_{1/2}$, $2p_{3/2}$,~\ldots, with the value of $j$ written as a subscript on the symbol specifying the location of the hole. Special symbols are customary, however, with the following correspondence:
\[
\begin{array}{ccccccccc}
1s_{1/2}&2s_{1/2}&2p_{1/2}&2p_{3/2}&3s_{1/2}&3p_{1/2}&3p_{3/2}&3d_{3/2}&3d_{5/2}\\
K&L_{\mathrm I}&L_{\mathrm{II}}&L_{\mathrm{III}}&M_{\mathrm I}&M_{\mathrm{II}}&M_{\mathrm{III}}&M_{\mathrm{IV}}&M_{\mathrm V}
\end{array}
\]
Levels with $n=4,5,6$ are denoted analogously by the letters $N,O,P$.

Levels having the same $n$ (and denoted by the same capital letter) lie close together and far from levels with other values of $n$. The reason is that inner electrons, being comparatively close to the nucleus, move in a field that is almost the unscreened nuclear field. Their states are therefore hydrogen-like, and \SourcePage{341}{344}to first approximation their energy is $-Z^2/2n^2$ (in atomic units), so it depends on $n$ alone.
